% Introduction — Paper 2, Medical Physics submission
% Humanised v1 (revised after the abstract/methods voice was settled).
% Three-paragraph structure mirroring Paper 1's intro layout.
% Formal-academic-active register, citation-bound.

Quantitative image-derived biomarker extraction has bifurcated into two paradigms whose relative merits the literature has only recently begun to compare directly. Handcrafted radiomics extracts hundreds of mathematically defined features --- first-order intensity statistics, three-dimensional shape descriptors, and second-order texture metrics computed over grey-level co-occurrence, run-length, size-zone, dependence, and neighbourhood matrices --- using Image Biomarker Standardisation Initiative (IBSI) 1.0~\cite{zwanenburg2020} reference implementations such as PyRadiomics~\cite{vangriethuysen2017}. Pretrained foundation-model (FM) embeddings replace the engineered feature dictionary with a dense vector produced by a vision encoder pretrained at scale on biomedical or natural images, including BiomedCLIP~\cite{zhang2024biomedclip}, RadImageNet-ResNet50~\cite{mei2022radimagenet}, the MedSAM image-encoder~\cite{ma2024medsam}, and a growing family of three-dimensional encoders~\cite{pai2024}. The implicit claim in much of the FM-radiomics literature is that FM embeddings offer both stronger discrimination and greater robustness to the upstream image-processing variability --- segmentation, scanner, reconstruction kernel --- that has driven the radiomics community to standardise around routine intraclass-correlation-coefficient (ICC) filtering at ICC $\geq 0.75$ before any downstream modelling~\cite{pavic2018,vantimmeren2020}.

Whether the FM stability claim holds for the specific perturbation mode that radiomics ICC-filtering exists to address --- segmentation-mask variability --- is unresolved in the peer-reviewed literature. The peer-reviewed precedent is Pai~\textit{et al.}~2024~\cite{pai2024}, who introduced FMCIB and reported ICC = 0.984 for its feature-based predictions on the RIDER same-day repeat-scan dataset, with significantly higher stability than a supervised CNN baseline under fifty trials of three-dimensional Gaussian seed-point jitter ($\sigma^2 = 16$ voxels per dimension) on the LUNG1 cohort. Pai~\textit{et al.}~2024 established that a well-pretrained 3D FM can be more stable than a supervised CNN under scanner-repeat and seed-point variability. They did not, however, perturb the segmentation mask itself, did not include an IBSI-aligned handcrafted-radiomics comparator on the same lesions, did not report per-dimension ICC across the embedding, and did not report pre-versus-post stability-filter downstream task performance. A contemporaneous preprint by the same group extends the 2024 benchmark to ten 3D FMs using global cosine similarity; this work has not undergone peer review and is referenced here only for context. Cosma~\textit{et al.}~2025~\cite{cosma2025}, working independently on breast MRI for triple-negative breast cancer subtype prediction, reported the counterintuitive finding that ICC-based filtering of handcrafted radiomic features can exclude features with predictive value --- raising the possibility that the inherited ICC-filter convention may be unnecessary when modern regularised classifiers are used. The central question of the present work, unanswered in the peer-reviewed literature, is therefore: how do FM embeddings behave when the segmentation mask itself is perturbed, and does ICC-based stability filtering produce a measurable downstream-task benefit when the same perturbation is applied to FM embeddings and to handcrafted radiomic features extracted from the same lesions?

We address this question on the LIDC-IDRI lung-CT cohort~\cite{armato2011}, which uniquely provides up to four independent expert segmentations per nodule. For 310 balanced nodules selected for clear binary malignancy labels, we extract IBSI-aligned PyRadiomics features under matched two-dimensional (central-slice) and three-dimensional (full-volume) configurations, and BiomedCLIP embeddings on the same central axial slice. For each nodule we apply four natural inter-reader perturbations together with seven controlled morphological perturbations and one Dice-calibrated synthetic mask anchored to that nodule's own empirical reader variability. We test four pre-registered hypotheses: that FM embeddings exhibit a higher proportion of stable features than radiomics under perturbation (H1); that downstream AUC degradation under ICC-based stability filtering is smaller for FM embeddings (H2); that post-filter AUC is non-inferior to pre-filter AUC at margin 0.020 (H3); and that the highest-ICC FM dimensions concentrate a disproportionate share of the malignancy-classification weight (H4). We complement the pre-registered tests with a Kolmogorov--Smirnov comparison of ICC distribution shapes (pre-specified before any analysis run) and a diameter-stratified subgroup analysis as the pre-specified mitigation of the LIDC size-malignancy confound. The findings reframe several common assumptions about FM-versus-radiomics stability and have direct methodological implications for the design of future foundation-model-versus-radiomics benchmarking studies.
